\section{A resolution of the local face module}

In this section, we prove Theorem~\ref{thm:resolution}, giving an explicit resolution of the local face module $L(\Gamma,E)$ by a subcomplex of the Koszul resolution of $k[\lk_{\Gamma}(E)]/(\theta_1, \dotsc, \theta_{d})$.  We continue to use the notation established above.  In particular, $\sigma\colon \Gamma \to 2^V$ is a quasi-geometric homology triangulation of the simplex with vertex set $V = \{v_1, \ldots, v_n\}$. We consider a face $E \in \Gamma$ with $d = n - |E|$ and $b = n - |\sigma(E)|$.  After reordering, we assume $\sigma(E)^c = \{v_{1}, \dotsc, v_{b}\}$. For $S \subset \{v_1, \ldots, v_d\}$, 
we consider the ideal $I_S \subset k[\lk_\Gamma(E)]$ given by
\begin{equation*}
I_S  \coloneqq  ( x^F : \, \sigma(F \sqcup E)^c \subset S).  
\end{equation*}

Let $\theta_1, \ldots \theta_d$ be a special l.s.o.p. for $k[\lk_\Gamma(E)]$. We may assume that 
\[
\supp(\theta_i) \subset \{ w \in \lk_\Gamma(E) : v_i \in \sigma(w) \},
\]
for $1 \leq i \leq b$.  For any $v_i \in S$, multiplication by $\theta_i$ gives a map $\lambda_i \colon I_{S} \to I_{S \smallsetminus \{v_i\}}$, and we consider the complex of graded $k[\lk_{\Gamma}(E)]$-modules
 \begin{equation} \label{eq:resolution}
 0  \to k[\lk_{\Gamma}(E)][-d] \to \bigoplus_{ |S|  = d - 1 } I_{S}[-(d-1)] \to \dotsb  \to  I \to L(\Gamma,E) \to 0,
\end{equation}
in which
 the differential restricted to $I_S$, for $S= \{v_{i_0}, \ldots, v_{i_k}\}$, with $i_0 < \cdots < i_k$, is $\oplus_{j = 0}^k (-1)^j \lambda_{i_j}$. 


\begin{example}
If $E$ is an interior face of $\Gamma$ then every l.s.o.p. is special, $I_S = k[\lk_{\Gamma}(E)]$ for all $S$, and \eqref{eq:resolution} is the Koszul resolution of $L(\Gamma,E) = k[\lk_{\Gamma}(E)]/(\theta_1,\ldots,\theta_{d})$.
\end{example}

\begin{proof}[Proof of Theorem \ref{thm:resolution}]
We must show \eqref{eq:resolution} is exact.
 We begin by considering two complexes of $k[\lk_{\Gamma}(E)]$-modules
 studied by Stanley and Athanasiadis. Recall that, for $U \subset V$, we write $\Gamma_U  \coloneqq  \sigma^{-1}(2^U)$.
 
Say $U \supset \sigma(E)$ and $U \smallsetminus \sigma(E) = \{v_{i_0}, \ldots, v_{i_k}\}$, with $i_0 < \cdots < i_k$. For $0 \leq j \leq k$, let $\rho_j \colon k[\lk_{\Gamma_U}(E)] \to k[\lk_{\Gamma_{U \smallsetminus \{v_{i_j}\}}}(E)]$ be the restriction map.  The first complex we consider is 
\small
\begin{equation}\label{eq:modcomp}
\begin{tikzcd}[column sep = tiny]
 k[\lk_{\Gamma}(E)]  \arrow[r] & \bigoplus\limits_{\substack{U \supset \sigma(E)  \\ \vert U \vert = n - 1}} k[\lk_{\Gamma_U}(E)] \arrow[r]  &\bigoplus\limits_{\substack{U \supset \sigma(E)  \\ \vert U \vert = n - 2}}  k[\lk_{\Gamma_U}(E)] \arrow[r]  &\cdots \arrow[r] & k[\lk_{\Gamma_{\sigma(E)}}(E)] \arrow[r] & 0,
\end{tikzcd}
\end{equation}
\normalsize
in which the differential restricted to $k[\lk_{\Gamma_U}(E)]$ is $\bigoplus_j (-1)^j \rho_j$.  Next, we consider its quotient by $(\theta_1,\ldots,\theta_{d})$:

\begin{equation}\label{eq:quotientedcomplex}
\begin{tikzcd}[column sep = small]
 \frac{k[\lk_{\Gamma}(E)]}{(\theta_1, \dotsc, \theta_{d})} \arrow[r] & \bigoplus\limits_{\mathclap{\substack{U \supset \sigma(E)  \\ \vert U \vert = n-1}}} \frac{k[\lk_{\Gamma_{U}}(E)]}{(\theta_1, \dotsc, \theta_{d})} \arrow[r]  & \bigoplus\limits_{\mathcal{\substack{U \supset \sigma(E)  \\ \vert U \vert = n-2}}} \frac{k[\lk_{\Gamma_{U}}(E)]}{(\theta_1, \dotsc, \theta_{d})} \arrow[r] &\cdots \arrow[r] & \frac{k[\lk_{\Gamma_{\sigma(E)}}(E)]}{(\theta_1, \dotsc, \theta_{d})} \arrow[r]  &0.
 \end{tikzcd}
\end{equation}

For any $U \subset V$, with $U \supset \sigma(E)$, let $S_U$ be defined as
\[
S_U  \coloneqq  (U \cap \{v_1, \ldots, v_b\}) \cup \{ v_{b+1}, \ldots, v_{d}\}.
\]
Then $\dim k[\lk_{\Gamma_U}(E)] = |S_U|$ and it follows that the restriction of $\theta_i$ to $\lk_{\Gamma_U}(E)$ is nonzero if and only if $v_i \in S_U$. Furthermore, $\{ \theta_i|_{\lk_{\Gamma_U}(E)} : v_i \in S_U \}$ is a special l.s.o.p. for $k[\lk_{\Gamma_{U}}(E)]$.  Stanley and Athanasiadis proved that both \eqref{eq:modcomp} and \eqref{eq:quotientedcomplex} are exact, and the kernel of the first arrow in \eqref{eq:quotientedcomplex} is $L(\Gamma,E)$. (We will recall the proofs below.) Using the additivity of Hilbert functions in exact sequences, they deduced that the Hilbert function of $L(\Gamma, E)$ satisfies \eqref{eq:localh} \cite{Stanley92, Athanasiadis12}.

With the goal of proving that \eqref{eq:resolution} is exact, we take Koszul resolutions of each term in \eqref{eq:quotientedcomplex} to build a double complex of $k[\lk_{\Gamma}(E)]$-modules. Since $k[\lk_{\Gamma_{U}}(E)]$ is Cohen-Macauley, the special l.s.o.p.  $\{ \theta_i|_{\lk_{\Gamma_U}(E)} : v_i \in S_U \}$ is a regular sequence. Hence the corresponding Koszul  complex $K^{\bullet}_U$
\begin{center}
\begin{tikzcd}[column sep = small]
0 \arrow[r] &k[\lk_{\Gamma_U}(E)]_{S_U} \arrow[r] &\bigoplus\limits_{\mathclap{\substack{S \subset S_U \\ \vert S \vert = \vert S_U \vert - 1}}} k[\lk_{\Gamma_U}(E)]_S \arrow[r]& \cdots \arrow[r]  &k[\lk_{\Gamma_U}(E)] \arrow[r]& \frac{k[\lk_{\Gamma_U}(E)]}{(\theta_1, \dotsc, \theta_{d})}  \arrow[r] &0,
\end{tikzcd}
\end{center}
is exact.  Here, for a graded module $M$ and a finite set $S$, we write $M_{S}  \coloneqq  M[-|S|]$. Replacing each term in \eqref{eq:quotientedcomplex} with its corresponding Koszul resolution, gives a complex of complexes
\begin{equation}\label{eq:koszul}
\begin{tikzcd}
K_{V}^{\bullet} \arrow[r] & \bigoplus\limits_{\mathclap{\substack{U \supset \sigma(E) \\ \vert U \vert = n - 1}}} K_U^{\bullet} \arrow[r] & \bigoplus\limits_{\mathclap{\substack{U \supset \sigma(E) \\ \vert U \vert = n - 2}}} K_U^{\bullet} \arrow[r] & \cdots \arrow[r] & K_{\sigma(E)}^{\bullet} \arrow[r] & 0,
\end{tikzcd}
\end{equation}
which may be expanded as the commuting double complex shown in Figure~\ref{fig:doublecx}.
\begin{figure}[h!] 
\adjustbox{scale=.8,center}{
 \begin{tikzcd}[column sep = tiny]
 0  &0  &0 & \cdots  & 0  \\
 \frac{k[\lk_{\Gamma}(E)]}{(\theta_1, \dotsc, \theta_{d})} \arrow[r] \arrow[u] & \bigoplus\limits_{\mathclap{\substack{U \supset \sigma(E)  \\ \vert U \vert = n-1}}} \frac{k[\lk_{\Gamma_{U}}(E)]}{(\theta_1, \dotsc, \theta_{d})} \arrow[r] \arrow[u] & \bigoplus\limits_{\mathcal{\substack{U \supset \sigma(E)  \\ \vert U \vert = n-2}}} \frac{k[\lk_{\Gamma_{U}}(E)]}{(\theta_1, \dotsc, \theta_{d})} \arrow[r] \arrow[u] &\cdots \arrow[r] & \frac{k[\lk_{\Gamma_{\sigma(E)}}(E)]}{(\theta_1, \dotsc, \theta_{d})} \arrow[r] \arrow[u] &0 \\
 k[\lk_{\Gamma}(E)]  \arrow[r] \arrow[u] & \bigoplus\limits_{\substack{U \supset \sigma(E)  \\ \vert U \vert = n - 1}} k[\lk_{\Gamma_U}(E)] \arrow[r] \arrow[u] &\bigoplus\limits_{\substack{U \supset \sigma(E)  \\ \vert U \vert = n - 2}}  k[\lk_{\Gamma_U}(E)] \arrow[r] \arrow[u] &\cdots \arrow[r] & k[\lk_{\Gamma_{\sigma(E)}}(E)] \arrow[r] \arrow[u] & 0 \\
\bigoplus\limits_{\mathclap{\substack{\vert S \vert = 1}}}k[\lk_{\Gamma}(E)]_S \arrow[r] \arrow[u] & \bigoplus\limits_{\substack{U \supset \sigma(E)  \\ \vert U \vert = n - 1}} \quad  \bigoplus\limits_{\mathclap{\substack{S \subset S_U \\ \vert S \vert = 1}}} k[\lk_{\Gamma_U}(E)]_S \arrow[r] \arrow[u] &\bigoplus\limits_{\substack{U \supset \sigma(E)  \\ \vert U \vert = n - 2}} \quad  \bigoplus\limits_{\mathclap{\substack{S \subset S_U \\ \vert S \vert = 1}}} k[\lk_{\Gamma_U}(E)]_S \arrow[r] \arrow[u] &\cdots \arrow[r] & \bigoplus\limits_{\mathclap{\substack{S \subset S_{\sigma(E)} \\ \vert S \vert =  1}}} k[\lk_{\Gamma_{\sigma(E)}}(E)]_S \arrow[r] \arrow[u] & 0 \\
\bigoplus\limits_{\mathclap{\substack{\vert S \vert = 2}}}k[\lk_{\Gamma}(E)]_S \arrow[r] \arrow[u] & \bigoplus\limits_{\substack{U \supset \sigma(E)  \\ \vert U \vert = n - 1}} \quad  \bigoplus\limits_{\mathclap{\substack{S \subset S_U \\ \vert S \vert = 2}}} k[\lk_{\Gamma_U}(E)]_S \arrow[r] \arrow[u] & \bigoplus\limits_{\substack{U \supset \sigma(E)  \\ \vert U \vert = n - 2}} \quad  \bigoplus\limits_{\mathclap{\substack{S \subset S_U \\ \vert S \vert = 2}}} k[\lk_{\Gamma_U}(E)]_S \arrow[r] \arrow[u] &\cdots \arrow[r] & \bigoplus\limits_{\mathclap{\substack{S \subset S_{\sigma(E)} \\ \vert S \vert = 2}}} k[\lk_{\Gamma_{\sigma(E)}}(E)]_S \arrow[r] \arrow[u] & 0. \\
\vdots  \arrow[u] & \vdots \arrow[u]  &\vdots \arrow[u] &\cdots  & \vdots \arrow[u]  \\
\bigoplus\limits_{\mathclap{\substack{\vert S \vert = d -1}}}k[\lk_{\Gamma}(E)]_S \arrow[r] \arrow[u] & \bigoplus\limits_{\mathclap{\substack{U \supset \sigma(E) \\ \vert U \vert = n - 1}}}k[\lk_{\Gamma}(E)]_{S_U} \arrow[u]\arrow[r] &0 \arrow[u] \\
k[\lk_{\Gamma}(E)]_{\{v_1, \ldots, v_d\}} \arrow[u] \arrow[r] &0 \arrow[u] \\
0 \arrow[u]
\end{tikzcd}}
\vspace{-15 pt}
\caption{The double complex obtained by taking the Koszul resolution of \eqref{eq:quotientedcomplex}.} \label{fig:doublecx}
\end{figure}
The columns of this complex are exact by construction. We claim that the rows are also exact, and prove this using ideas from \cite[Theorem~4.6]{Stanley92}. First, we show that all rows except for the top row are exact. Choose a subset $S$ of $\{v_1, \ldots, v_d\}$, and consider the piece of the complex indexed by $S$:
\begin{equation}\label{eq:generalS}
\begin{tikzcd}[column sep = small]
 k[\lk_{\Gamma}(E)]_S  \arrow[r]  & \bigoplus\limits_{\substack{ S \subset S_U   \\ \vert U \vert = n - 1}} k[\lk_{\Gamma_U}(E)]_S \arrow[r]  &\bigoplus\limits_{\substack{S \subset S_U \\ \vert U \vert = n - 2}}  k[\lk_{\Gamma_U}(E)]_S \arrow[r]  &\cdots 
 \arrow[r]  & 0.
\end{tikzcd}
\end{equation}



When $S = \emptyset$, we obtain (\ref{eq:modcomp}). Observe that the complex (\ref{eq:generalS}) is multigraded by $\mathbb{N}^m$, where $m$ is the number of vertices of $\lk_{\Gamma}(E)$. Explicitly, $\deg x_1^{\alpha_1} \dotsb x_m^{\alpha_m} = (\alpha_1, \dotsc, \alpha_m)$. Therefore it suffices to show exactness on graded pieces. Fix $\alpha = (\alpha_1, \dotsc, \alpha_m)$. By the definition of the face ring, every term of (\ref{eq:generalS}) will have $0$ in the graded piece corresponding to $\alpha$ unless the set of vertices with $\alpha_i \not= 0$ forms a face $F$, in which case the $\alpha$-graded part can be identified with the augmented cochain complex of a simplex, indexed by all $U$ that contain $\sigma(E) \cup \sigma(F) \cup S$, and hence is exact.

We now recall the proof that the top row of the double complex, \eqref{eq:quotientedcomplex}, is exact.
\begin{center}
\begin{tikzcd}[column sep = small]
 \frac{k[\lk_{\Gamma}(E)]}{(\theta_1, \dotsc, \theta_{d})} \arrow[r]  & \bigoplus\limits_{\mathclap{\substack{U \supset \sigma(E)  \\ \vert U \vert = n-1}}} \frac{k[\lk_{\Gamma_{U}}(E)]}{(\theta_1, \dotsc, \theta_{d})} \arrow[r] & \bigoplus\limits_{\mathcal{\substack{U \supset \sigma(E)  \\ \vert U \vert = n-2}}} \frac{k[\lk_{\Gamma_{U}}(E)]}{(\theta_1, \dotsc, \theta_{d})} \arrow[r]  &\cdots \arrow[r] & \frac{k[\lk_{\Gamma_{\sigma(E)}}(E)]}{(\theta_1, \dotsc, \theta_{d})} \arrow[r]  &0
\end{tikzcd}
\end{center}
The proof involves showing that the quotients of (\ref{eq:modcomp}) by $(\theta_{d}, \dotsc, \theta_{d - (r-1)})$ is exact by induction on $r$. The case of $r = 0$ is the exactness of the second row. 

Now assume that (\ref{eq:modcomp}) remains exact after quotienting by $(\theta_{d}, \dotsc, \theta_{d - (r-1)})$. Let $C^i$ denote the $i$th term of (\ref{eq:modcomp}) tensored with $k[\lk_{\Gamma}(E)]/(\theta_{d}, \dotsc, \theta_{d - (r - 1)})$. By the induction hypothesis, we have an exact sequence
\begin{equation*}
C^{\bullet}\colon \enskip  C^0 \to C^1 \to \dotsb \to C^{b} \to 0.
\end{equation*} 
Set $m = d - r$. Recall that 
$\theta_i = 0 \in k[\lk_{\Gamma_{U}}(E)]$ if $v_i \notin S_U$, and that $\{ \theta_i|_{\lk_{\Gamma_U}(E)} : v_i \in S_U \}$ is a special l.s.o.p. for $k[\lk_{\Gamma_{U}}(E)]$. 
Also, for $\sigma(E) \subset U$, $v_m \notin S_U$ if and only if $v_m \notin U$. Hence, we have an exact sequence

\begin{equation} \label{eq:multbym}
0 \to B^{\bullet} \to C^{\bullet} \xrightarrow{\theta_{m}} C^{\bullet} \to C^{\bullet}/(\theta_{m}) \to 0,
\end{equation}
where 
\[
B^{i} = \bigoplus_{\substack{U \supset \sigma(E), \enskip \vert U \vert = n-i \\ v_m \not \in U}}k[\lk_{\Gamma_U}(E)]/(\theta_{d}, \dotsc, \theta_{m + 1}).
\]
For example, when $m > b$, $v_m \in \sigma(E)$ and $B^{\bullet} = 0$.
Up to signs and a degree shift, we can then identify $ B^{\bullet}$ with the complex (\ref{eq:modcomp}) for $\Gamma|_{\{v_m\}^c}$ quotiented by $(\theta_{d}, \dotsc, \theta_{m + 1})$. Then $ B^{\bullet}$ is exact by the induction hypothesis applied to $\Gamma|_{\{v_m\}^c}$. By breaking (\ref{eq:multbym}) up into two short exact sequences we see that $H^i(C^{\bullet}/(\theta_{m})) \cong H^{i + 2}(B^{\bullet}) = 0$ as desired.

Now that we know the exactness of (\ref{eq:koszul}), let 
\begin{equation*} \begin{split}
A^{\bullet} &= \ker \Bigg( K_{V}^{\bullet} \to \bigoplus\limits_{\mathclap{\substack{U \supset \sigma(E) \\ \vert U \vert = n - 1}}} K_U^{\bullet} \Bigg).
\end{split}\end{equation*}
Then, by construction, we have an exact sequence of complexes 
\begin{equation*}
\begin{tikzcd}
0 \arrow[r] & A^{\bullet} \arrow[r] &K_{V}^{\bullet} \arrow[r] & \bigoplus\limits_{\mathclap{\substack{U \supset \sigma(E) \\ \vert U \vert = n - 1}}} K_U^{\bullet} \arrow[r] & \bigoplus\limits_{\mathclap{\substack{U \supset \sigma(E) \\ \vert U \vert = n - 2}}} K_U^{\bullet} \arrow[r] & \cdots \arrow[r] & K_{\sigma(E)}^{\bullet} \arrow[r] & 0.
\end{tikzcd}
\end{equation*}
As above, we repeatedly apply the long exact sequence on cohomology to see that $A^{\bullet}$ is exact. 
We may then identify $A^{\bullet}$ with the exact sequence
$$0  \to k[\lk_{\Gamma}(E)][-n] \to \oplus_{ |S|  = d - 1 } I_{S}[-(n-1)] \to \dotsb  \to \oplus_{|S| = 1} I_{S}[-1] \to I \to A^0 \to 0.$$ Since $I$ surjects onto $A^0$ and $A^0 \subset k[\lk_{\Gamma}(E)]/(\theta_1, \dotsc, \theta_{d})$, we conclude that $A^0 = L(\Gamma,E)$, as required.
\end{proof}
